\documentclass[10pt]{article}
\usepackage[margin=20mm]{geometry}
\usepackage{amsmath,amssymb,mathtools,fontspec,unicode-math,xurl,hyperref}
\IfFontExistsTF{DejaVu Serif}{\setmainfont{DejaVu Serif}}{\setmainfont{Latin Modern Roman}}
\setmathfont{latinmodern-math.otf}
\setlength{\parindent}{0pt}\setlength{\parskip}{6pt}\sloppy
\title{Initial scientific review
The master equation for mean field game systems with fractional and nonlocal diffusions}
\author{Assigned initial-review executor}
\date{}
\begin{document}
\maketitle
Paper. Espen R. Jakobsen and Artur Rutkowski, DOI 10.4171/JEMS/1681. Exact inspected source: arXiv:2305.18867v2, 24 January 2025, 61 pages. The publisher records online publication on 30 July 2025; equivalence with its subscription version has not been established.\par
\section*{Assessment and scope}
The paper develops a substantial analytical route from whole-space nonlocal MFG systems to a classical master equation without moment assumptions. Three Major issues concern a tightness estimate, the value-dependent uniqueness argument, and the literal quantifiers in the diffusion assumption. Further corrections are required in several auxiliary results. No pair of distinct solutions to the master equation is exhibited. The main theorem is therefore not declared false on the evidence below.\par
All 61 pages, including Appendices A–D and references, were read as extracted text. The earlier EMFILE failure and incomplete state are preserved. Scientific decision remains pending; this is an initial review, not independent acceptance. No author code was executed.\par
\section*{Coverage of the scientific argument}
The analysis concerns constant-coefficient Lévy generators with a positive semidefinite local part and the stated Lévy integrability condition. The metric on probability measures is the bounded Lipschitz metric \(d_0\), not a Wasserstein metric imposing moments. The subadditive smooth tightness function replaces polynomial moments. All quantifiers in (H1)–(H3), (F1), (F2), (G1), (G2) and (M1)–(M4) were checked. In particular, Sections 5–6 require a Hamiltonian independent of the value variable, uniform convexity in momentum, and the stronger dual-space positivity (M2). Lasry–Lions monotonicity alone is not silently substituted for (M2). The measure derivative has a gauge freedom; the authors’ discussion of normalization and the odd-kernel example was checked.\par
Section 2 distinguishes the full dual \(C_b^{-\gamma}\) from the measure-representable subspace \(\mathfrak C_b^{-n}\). A functional at infinity need not mollify to an \(L^1\) function. The initial datum in Theorem 5.5 is explicitly measure-representable; a proposed objection using a Banach limit was rejected for violating this hypothesis. Iterated spatial Hölder norms, weak convergence, and differentiation along convex combinations of measures were followed with their specified domains.\par
Section 3 uses integrability of \(t^{-1/\alpha}\), hence \(\alpha>1\), for fixed points and Duhamel estimates. The positive, source-free probability equation supports the moment-free Schauder construction in Theorem 4.3. This must be distinguished from the signed equation with an additive divergence source in Lemma 3.5, where P1 applies. Uniform spatial regularity of the Hamilton–Jacobi equation and its time continuity in a slightly weaker norm were traced separately from the overstrong time-derivative modulus in P6.\par
In Section 4 the existence map has uniform drift, tightness and \(d_0\) equicontinuity bounds. For uniqueness, the last inference of Lemma 4.2 is unsupported as detailed in P2. For a Hamiltonian independent of the value variable, the usual monotonicity identity does give the energy estimate used later; that case does not need the defective extension.\par
Theorem 5.3 combines the backward estimate with a forward duality estimate containing the square root of an energy term, and absorbs the resulting product. Theorem 5.5 uses a regular-data fixed point, dual backward test functions, (M2), and ellipticity of \(\Gamma\). The energy/absorption calculation (5.15)–(5.21) was checked. Its compactness step on page 37 expressly invokes Lemma 3.5(e), so P1 is a core dependency. Approximation of initial functionals and sources uses the representable subspace and weaker dual norms; Appendix A needs the correction in P4.\par
Lemma 5.7’s weighted smoothing requires \(\gamma<\alpha\); the beta-integral at the initial time and the divergence-source term were checked. The passage removing an arbitrarily small loss on page 41 needs uniform constants or a direct estimate; no counterexample to the intended smoothing result is asserted. The time origin in Corollary 5.8 must be shifted as in P7. Dirac difference quotients in Proposition 5.10 give the first and second \(y\) derivatives, and the integrable smoothing exponent \(1+\sigma<\alpha\) controls its coefficient differences. Proposition 5.11 needs a correct composition estimate, P5, to obtain its superlinear measure Taylor remainder.\par
Section 6 differentiates the MFG solution along its measure flow. The weak Fokker–Planck formulation and the bounded, uniformly continuous derivatives of the candidate justify the limiting integrals. The uniqueness construction solves the induced Fokker–Planck equation for a competing classical solution and returns to the value-independent MFG uniqueness. Appendices B–D were read: the weaker spatial-norm time modulus is valid in the displayed Duhamel argument; mild-to-weak conversion requires the stated test-function regularity and integrable kernel singularities. The initial-time stability argument uses the spatial modulus, not the false forcing-independent modulus for \(\partial_tu\).\par
\section*{Findings with decisive evidence}
\subsection*{P1. Major: a divergence source has uncontrolled weighted absolute mass}
Location. Lemma 3.5(e), page 24; its proof, page 26; Theorem 5.5, compactness step 2, page 37.\par
The proposed bound controls \(\int\psi|\rho(t)|\) using only the initial weighted mass, unweighted \(L^1\) bounds for \(V_2,\rho\), and fixed coefficients. Translating a source disproves that bound in the Gaussian example expressly included in the paper. Here and in P6, Gaussian tests use the necessary per-order interpretation of (K) explained in P3; they are not claimed to satisfy the impossible literal all-order hypothesis.\par
Set \(d=1\), \(T=1\), \(\mathcal L=\Delta\), \(V_1=0\), \(\rho_0=0\), and \(V_2(t,x)=G_1(x-R)\), where \(G_s(x)=(4\pi s)^{-1/2}e^{-x^2/(4s)}\). These data are smooth and \(\left\|V_2(t)\right\|_1=1\). The solution of (3.10) is \[\rho_R(t,x)=\int_0^t\partial_xG_{1+s}(x-R)\,ds.\] Put \(h(y)=\int_0^1\partial_yG_{1+s}(y)\,ds
=\tfrac12[\operatorname{erf}(y/(2\sqrt2))
-\operatorname{erf}(y/2)]\). Its sign is opposite to \(y\), and direct integration gives \[M:=\int_\mathbb R|h(y)|\,dy=\frac{2(\sqrt2-1)}{\sqrt\pi}>0,
 \qquad \int_\mathbb R|y|\,|h(y)|\,dy=1 .\] For \(\psi(x)=\sqrt{1+x^2}\), which is smooth, nonnegative, subadditive and has bounded first two derivatives, \[\int_\mathbb R\psi(x)|\rho_R(1,x)|\,dx\ge RM-1\longrightarrow\infty .\] Every quantity permitted on the right of Lemma 3.5(e) remains bounded independently of \(R\), including \(\sup_{t\le1}\left\|\rho_R(t)\right\|_1\le M\). The initial weighted mass is zero and the jump measure is zero. If a probability initial datum is preferred, add \(G_1\) initially; the solution acquires \(G_{t+1}\) and the same lower bound loses only a fixed finite weighted mass.\par
Testing a signed equation with \(\psi\) estimates a signed moment, not its absolute moment. The cited positive, source-free argument does not establish the asserted estimate. The compactness input therefore requires replacement. For the actual fixed-point family, the special source \(m\Gamma Dz+c\) has more structure: uniform tails of the fixed \(m,c\) and bounded \(Dz\) can support a Duhamel tail argument. Such an argument must be supplied; the counterexample does not rule it out or refute Theorem 6.2.\par
\subsection*{P2. Major: the last Gronwall inference does not follow}
Location. Last display of page 31 and first display of page 32, Lemma 4.2; Theorem 4.3(b).\par
Writing \(A(s)=\int u^+m^++\int u^-m^-\ge0\), the proof obtains \[A(s)+c_2\int_s^T A(t)\,dt\ge J(s)\] and concludes \(J(t_0)\le C A(t_0)\). This inference is false even with \(J\) nonnegative, decreasing, absolutely continuous and \(J(T)=0\), as in the proof. Take \(t_0=0\), \(T=c_2=1\), \[A(s)=s^2,\qquad J(s)=\int_s^1t^2\,dt=(1-s^3)/3 .\] The premise holds for every \(s\), while \(A(0)=0<J(0)=1/3\). Thus Gronwall cannot justify the conclusion. This is a counterexample to the actual scalar inference, not a constructed nonunique MFG solution.\par
The general value-dependent uniqueness conclusion under (M4) needs an additional estimate or a different argument. For \(H_2(x,u)=\beta u\) with constant \(\beta\), an integrating factor can handle the corresponding energy identity; this does not prove the stated general \(x,u\)-dependent case. The value-independent uniqueness used in Sections 5–6 has a separate valid monotonicity argument.\par
\subsection*{P3. Major as written: the diffusion hypothesis has an empty literal scope}
Location. Assumption (K), Section 2.5.2, page 14, and the examples on page 16.\par
The stated quantifier is one finite \(\mathcal K\) for every multi-index \(\beta\). A Lévy characteristic function has the form \(\widehat K_t(\xi)=e^{-t\Psi(\xi)}\), which is nonzero at each finite \(\xi\). The proposed \(L^1\) estimate implies, for all \(n\), \[(|\xi_1|t^{1/\alpha})^n|\widehat K_t(\xi)|\le\mathcal K .\] Choosing \(|\xi_1|>t^{-1/\alpha}\) gives a contradiction as \(n\to\infty\). In particular even the Gaussian example fails the literal hypothesis. This is an exact Fourier argument, not a numerical failure to find an operator.\par
The intended formulation is presumably a constant \(\mathcal K_\beta\) for each derivative order, or a common constant for a specified finite range of orders. That change must be explicit wherever theorem hypotheses refer to (K). For mixed anisotropic orders, estimates using the smallest order also need a finite time horizon or appropriate large-time bounds. The Major classification concerns the actual scope of the stated main theorem; it does not claim the intended finite-order semigroup theory is impossible.\par
\subsection*{P4. Minor norm correction, with a replacement for the closure argument}
Location. Lemma 2.6, page 8; Appendix A, pages 50–51, especially the inference after (A.6).\par
With the paper’s norm \(\left\|\phi\right\|_{C_b^1}=\left\|\phi\right\|_\infty+\left\|\phi'\right\|_\infty\), the functional \(\rho(\phi)=\phi(0)+\phi'(0)\) has norm one. Every representation \(\rho(\phi)=\int\phi\,d\mu_0+\int\phi'\,d\mu_1\) has \(\left\|\mu_0\right\|_{TV}\ge1\) by testing the constant one. Taking \(\phi_\epsilon(x)=\epsilon q(x/\epsilon)\) with \(q\in C_c^\infty\), \(q(0)=0\), \(q'(0)=1\), and letting \(\epsilon\downarrow0\), gives \(\mu_1(\{0\})=1\). Hence \(\left\|\mu_1\right\|_{TV}\ge1\); the least sum is two, attained by \(\mu_0=\mu_1=\delta_0\). The claimed exact identity with the sum of total variations is false. The dual of the sum norm on the jet product uses a maximum norm.\par
In addition, equality of each chosen lift’s norm with its functional’s norm does not imply the lifts are Cauchy. The closure conclusion itself remains correct: restriction identifies the representable subspace isometrically with \((C_0^n)^*\). Hahn–Banach on the finite jet product and Riesz representation show surjectivity. For a representable functional, cutoffs \(\chi(x/R)\phi\) and dominated convergence give equality of its \(C_b^n\) and restricted \(C_0^n\) norms, since the cutoff derivative error tends to zero. Thus this subspace is complete and closed without choosing Cauchy representatives. This replacement preserves the closure used downstream, while the exact norm formula needs correction.\par
\subsection*{P5. Minor: the composition estimate omits individual Lipschitz bounds}
Location. Lemma 5.9, page 42; applications on pages 44 and 46.\par
Choose \(\sigma=1/2\), \(\epsilon=1/4\), \(f(x,y)=\sin y\), \(y_1(x)=\sin(Nx)\) and \(y_2(x)=y_1(x)+\pi\). All functions have the stated bounded smooth regularity. The norm of \(y_1-y_2\) is the constant \(\pi\), and \(\left\|f\right\|_{C_b^{3/4}}\le1+2^{1/4}\). The left side of the proposed estimate is at least \[4\sin(1)(N/\pi)^{1/2},\] by testing \(x=\pi/(2N)\) and \(3\pi/(2N)\), whereas its right side is at most \(2(1+2^{1/4})(2+\sqrt\pi)\pi^{1/4}\), independently of \(N\). Already \(N=10^6\) contradicts the displayed inequality.\par
A correct composition/interpolation estimate includes the individual Lipschitz norms of \(y_1,y_2\). In the actual applications these are gradients of solutions with a uniform \(C_b^{3+\sigma}\) bound, so that dependence can be absorbed in the existing generic data constant. The counterexample refutes the standalone lemma; it does not by itself refute the superlinear remainder claimed in Proposition 5.11. Its proof and constants should be rewritten.\par
\subsection*{P6. Minor corrections to temporal and endpoint estimates}
Locations. Lemma 3.3(b), page 21; Theorem 3.10, (3.27), page 29; Lemma 3.5(c), page 25.\par
For \(\mathcal L=\Delta\), zero initial data and zero drift, \(f_n(t,x)=\sin(nt)\) gives \(z_n(t,x)=(1-\cos nt)/n\). The forcing’s spatial norms are one, but \(|\partial_tz_n(\pi/(2n))-\partial_tz_n(0)|=1\). There is no common vanishing time modulus for these derivatives depending only on the displayed spatial bounds. The same spatially constant example works with \(H(p)=|p|^2\). Each derivative is continuous; the error is the asserted uniformity. The forcing’s own time modulus must enter. The weaker spatial-norm modulus proved in Appendix B is not contradicted, and is the part used in Appendix D.\par
At \(\alpha=2\), the page-25 integral is \(\int_\epsilon^t r^{-1}\,dr=\log(t/\epsilon)\), not a constant times \(\epsilon^{\,0}\). The displayed argument establishes at best a logarithmic endpoint modulus, or any smaller positive Hölder exponent, which still suffices for translation compactness. No counterexample to the entire compactness conclusion is claimed.\par
\subsection*{P7. Minor: time origin and additional formula corrections}
Locations. Corollary 5.8, page 42; mild HJB formulas on pages 20 and 29; Lemma 2.3, page 7.\par
Corollary 5.8 needs \((t-t_0)^{\gamma/\alpha}\) in place of \(t^{\gamma/\alpha}\). For \(\mathcal L=\Delta\), \(F=G=0\), \(H(p)=|p|^2/2\) and the zero backward solution, the forward datum \(\rho_0=\partial_{xx}\delta_0\) is measure-representable of order two. At \(t_0>0\), set \(\tau=t-t_0\) and test \(\rho(t)=\partial_{xx}G_\tau\) with \(\phi_\tau(x)=\tau^{1/4}\cos(x/\sqrt\tau)\). Its \(C_b^{1/2}\) norm is uniformly bounded for \(0<\tau\le1\), whereas the pairing is \(-e^{-1}\tau^{-3/4}\). Thus the printed weight cannot stay bounded at \(t_0\). The intended integrability used later survives the shift.\par
For the forward-time equation \(\partial_tu-\mathcal Lu+H=f\), the Duhamel source is \(f-H\); the printed \(H-f\) has the opposite sign. With \(H(u)=u\), \(f=0\), \(u(0)=1\), it gives \(e^t\) instead of \(e^{-t}\). This corrects the displayed fixed-point map, not the norm-based contraction argument.\par
Finally the arbitrary rate \(\epsilon^\delta\) in Lemma 2.3 needs the standard restriction on the smoothing exponent. For an even nontrivial mollifier, \(\rho=\delta_0\), \(\gamma=0\) and the test \(\cos x\), the error is \(\frac12\epsilon^2\int y^2\eta(y)\,dy+o(\epsilon^2)\); it is not \(O(\epsilon^3)\). Convergence in the weaker dual norm still follows using a smaller positive exponent.\par
\section*{Evidence, limitations and disposition}
The independent program checks the Gaussian mass identities and translated weighted integrals, the scalar polynomial inference, the high-frequency composition bound, the temporal forcing example, and the Fourier growth. It is reviewer-written code; numerical quadrature supports exact analytical arguments and is not presented as a rigorous numerical enclosure or a simulation of an MFG solution.\par
The arXiv version record, exact-v2 HTML, publisher metadata and the author’s publication page were searched for supplements, code and corrections. No separate relevant supplement, code archive or erratum was located in that bounded search. Absence of code is not a mathematical defect. Reference [34]’s full proof and all external functional-analytic or semigroup results were not independently reproved; the report does not certify them. The prospective \(N\)-player convergence work is outside this paper’s established results.\par
The required revisions must separate the value-dependent uniqueness claim from the value-independent master equation, repair the signed-source compactness argument, and state a nonempty heat-kernel hypothesis. The remaining corrections have the narrower effects described above. Scientific decision remains pending; human reviews are empty. Exact source, TeX, compiler, dependencies, PDF, actual-view records and preserved unsuccessful visual-check records are bound in the per-paper provenance and inventory.\par
\end{document}
